\documentclass[11pt]{article}
\usepackage[margin=25mm]{geometry}
\usepackage{fontspec}
\setmainfont{Latin Modern Roman}
\usepackage{amsmath,amssymb,mathtools}
\usepackage{unicode-math}
\setmathfont{Latin Modern Math}
\usepackage{xurl,hyperref}
\hypersetup{colorlinks=true,urlcolor=blue}
\setcounter{secnumdepth}{0}
\setlength{\parindent}{0pt}
\setlength{\parskip}{5pt}
\setlength{\emergencystretch}{3em}
\newcommand{\sref}[2]{\hyperlink{source:#1:#2}{[S#2]}}
\newcommand{\cataloguescope}[1]{\paragraph{Recorded target.}#1}
\newcommand{\R}{\mathbb R}
\newcommand{\C}{\mathbb C}
\newcommand{\N}{\mathbb N}
\newcommand{\Z}{\mathbb Z}
\newcommand{\Q}{\mathbb Q}
\newcommand{\F}{\mathbb F}
\newcommand{\eps}{\varepsilon}
\begin{document}
\section{Interior Singularity Structure of Dynamically Formed Black Holes}\label{statement}

\subsection{Definitions and mathematical statement}

\paragraph{Shared notation.}$\mathbb N=\{1,2,\ldots\}$; the standard number fields have their usual meanings. All additional conventions are stated locally.


A Cauchy horizon is the boundary of the domain determined by initial data. A weakly singular horizon may admit a continuous metric extension while curvature or higher derivatives diverge. Gravitational collapse produces a black-hole interior whose future boundary can include both such null boundaries and stronger terminal singularities.

\paragraph{Catalogue wording retained for traceability.}
What is the global singularity structure inside a black hole formed by gravitational collapse, in particular the relationship between a weakly singular Cauchy horizon, stronger terminal singularities, and the center?

\paragraph{Scope and source qualification.}This is a structural research question. The matter model and regularity notion must be fixed before claiming a theorem for a particular collapse class.
\subsection{Short English statement}
Describe how weak Cauchy-horizon singularities, stronger singularities and the center fit together inside dynamically formed black holes.
\subsection{Research attempt}
\paragraph{Model, date and investigated scope.}
OpenAI GPT-6 Astra (configured model \texttt{gpt-6-astra}), ultra reasoning; generated 11 October 2026. Constructed a continuously extendible positive-radius warped metric with a divergent curvature component, clarifying distinct singularity regularities without a collapse model.

\paragraph{Sources and literature check.}
Read the supplied definition, catalogue wording, source qualifications and reference entries for this question. The historical source/context and inherited status are retained; no fresh literature or novelty certification is claimed. No external theorem is used without proof in this attempt; explicit model assumptions are distinguished from facts about the source model.

\paragraph{Separating metric continuity from curvature regularity.} A local warped-product diagnostic makes the distinction precise. For \(v>0\), let
\[
 g=-2\,du\,dv+r(v)^2d\Omega^2,\qquad r(v)=r_0+v^\alpha,
 \quad r_0>0,\quad0<\alpha<1.
\]
The metric extends continuously and nondegenerately to \(v=0\), and the symmetry-sphere area tends to \(4\pi r_0^2>0\). Its angular Christoffel component is \(\Gamma^\theta_{v\theta}=r'/r\). Substituting this into the curvature formula yields
\[
 R_{v\theta v\theta}=-rr''=-r\alpha(\alpha-1)v^{\alpha-2},
\]
which is unbounded in this coordinate frame that remains regular for the continuous extension. The derivative \(r'\) also diverges.
\paragraph{What it demonstrates.} A boundary at positive area radius may be continuous at the metric level and singular at curvature level; it need not be the center, where area radius vanishes. Hence a proof of one type of extendibility cannot by itself classify the global terminal boundary or identify its endpoints with the center.
\paragraph{Source comparison.} The displayed metric is a geometric diagnostic, not a verified Einstein--matter collapse solution or a Cauchy horizon in a globally specified development. The first gap is controlling the actual coupled field equations from regular collapse data across regions where different singularity types meet. Existence of a weak null segment, its relation to stronger terminal segments, and the global center geometry all remain. No curvature contraction or coordinate-invariant blowup theorem is claimed solely from this component calculation.

\paragraph{Verification and adversarial audit.}
Checked nondegenerate continuous extension, positive limiting area and the explicit curvature component; withheld Einstein-solution and global-Cauchy-horizon claims. This is a GPT-6 Astra Ultra self-audit, not independent human verification. Prompt A drafting followed by prompt B scope/dependency/boundary self-audit. Checked nondegenerate continuous extension, positive limiting area and the explicit curvature component; withheld Einstein-solution and global-Cauchy-horizon claims. The remaining gap is: Control the actual collapse equations globally to identify and connect weak null boundaries, terminal singularities and the center. Unresolved status and zero original-target completion are retained.

Coordinator checked the warped angular curvature component and nondegenerate continuous metric limit, with no unproved field equation or invariant blowup assertion. This additional check was model-assisted, not independent human verification.

\paragraph{Tools, process and contributions.}
Python 3 and local file-reading/writing tools for manuscript serialization and structural checks; mathematical calculations were symbolic reasoning, with no numerical experiment or theorem prover. Contribution prompt A was followed by prompt B. GPT-6 Astra with ultra reasoning drafted the research and performed its stated self-audit; any additional coordinator contribution is identified in the verification and audit. A GPT-6-based Codex coordinator assembled the manuscript and checked repository integration; its exact variant and reasoning configuration were not exposed. No human mathematical contribution, independent human proof check or formal proof certification is supplied. Complete derivations and the consulted source references are included here; private reasoning transcripts are not published.

\paragraph{Final status and first remaining gap.}
\textbf{UNRESOLVED}. Control the actual collapse equations globally to identify and connect weak null boundaries, terminal singularities and the center.

\section{Completion Estimate}
\noindent\textbf{COMPLETION: 0\%}
The proved calculation is an elementary auxiliary result or a test of a restricted strategy; no new portion of the full source target is established.

\subsection{Sources}
\begin{itemize}
\item[\hypertarget{source:statement:1}{S1}] ULAM.AI and the UnsolvedMath Contributors, UnsolvedMath, catalogue entry. Catalogue curation under CC BY 4.0; definitions and English exposition added locally.
\url{https://huggingface.co/datasets/ulamai/UnsolvedMath}.
\item[\hypertarget{source:statement:2}{S2}] Mathematical Aspects of General Relativity (2021), Oberwolfach Reports, DOI 10.4171/owr/2021/40, pp. 2157--2267. Original report and mathematical context.
\url{https://doi.org/10.4171/owr/2021/40}.
\end{itemize}
\label{end:statement}
\end{document}
